\section{Yang--Mills Mass Gap via Lattice $\to$ OS $\to$ Spectral Gap}
\label{sec:goertzel-ym-state}

This section is a research-program lens on Goertzel's route to the Yang--Mills
mass gap: build a lattice gauge theory, pass through Osterwalder--Schrader (OS)
reconstruction, and extract a spectral gap that survives the continuum limit. It
is written to direct creative effort. The route now has a genuine
\emph{finite-lattice landmark} --- a fully machine-checked spectral-gap lower
bound for a finite $\mathbb{Z}_2$ strong-coupling model --- and a sharply
identified open crux: continuum gap-uniformity, i.e.\ the asymptotic-freedom
wall. Every conjecture below is labelled. Lean identifiers cite the live
\texttt{QuantumTheory/YangMills} development.

\subsection{The creative bet}

The Clay statement, in spectral terms, is that the Hamiltonian has no spectrum in
an interval $(0,\Delta)$ above the vacuum. The route's bet is the standard
constructive-QFT strategy made fully constructive and auditable: (i)~define a
finite lattice gauge model with a reflection-positive action; (ii)~run OS
reconstruction to get a transfer operator and a Hilbert space; (iii)~prove a
\emph{transfer-matrix spectral gap} on the lattice; (iv)~prove that gap survives,
uniformly, as the lattice spacing $a\to 0$. The bet is that strong-coupling
expansions give step~(iii) cleanly, and that a renormalization-group (RG)
contraction controls step~(iv).

The genuinely new, machine-checked result is the finite-lattice endpoint of
steps (i)--(iii) for the abelian toy $\mathbb{Z}_2$:
\[
  \verb|z2StrongCoupling_latticeYangMills_massGap_landmark|.
\]
This packages an explicit $\mathbb{Z}_2$ strong-coupling model with a nontrivial
configuration, reflection positivity, an OS transfer operator, and a
dimensionless transfer gap $\ge 1$. Supporting it:
\begin{itemize}\itemsep1pt
  \item \verb|z2StrongCouplingGap_eq_one| and
    \verb|z2StrongCouplingGap_lower_bound_one| --- the gap equals $1$ and is
    bounded below by $1$ (with $\lambda_1=1$, $\lambda_2=e^{-1}$, gap
    $=-\log(\lambda_2/\lambda_1)=1$);
  \item \verb|z2StrongCouplingKernel_reflectionPositive| and the OS bundle
    \verb|FiniteOSReconstruction| (\verb|OSReflectionPositive|,
    \verb|OSSymmetric|, \verb|OSTransferSelfAdjoint|) --- the finite OS
    interface really holds for this model;
  \item the explicit canary \verb|z2StrongCouplingPositiveCanary| and proof node
    \verb|yangMillsZ2StrongCouplingGapNode|.
\end{itemize}
This is the first of the four routes to carry a positive, end-to-end finite
construction with a strictly positive gap. That is a real foothold.

\subsection{The OS-gate canaries (why we cannot overclaim)}

The development pairs the landmark with a battery of \emph{canaries} that block
premature promotion of a finite result to a continuum claim. Each canary is a
positive/negative pair showing that an interface is necessary but not
sufficient:
\begin{itemize}\itemsep1pt
  \item OS interfaces alone do not clear the continuum gate ---
    \verb|currentYangMillsOSGateCanaryAuditPacket|,
    \verb|toyYangMillsOSOnlyWorldInterfaces| (blocked) vs
    \verb|toyYangMillsCompletePromotionGate| (cleared);
  \item the continuum measure needs \emph{both} tightness and a limit ---
    \verb|currentYangMillsContinuumMeasureCanaryAuditPacket|, with sampler-only,
    tightness-without-limit, and limit-without-tightness all blocked;
  \item Euclidean covariance needs action \emph{and} preservation, and needs the
    continuum measure upstream ---
    \verb|currentYangMillsEuclideanCovarianceCanaryAuditPacket|;
  \item reflection positivity at the continuum needs functional \emph{and} cone
    interfaces, not a sample flag ---
    \verb|currentYangMillsReflectionPositivityCanaryAuditPacket|.
\end{itemize}
These canaries are the route's honesty mechanism: they enumerate, finitely,
everything the continuum endpoint still requires.

\subsection{The open crux: continuum gap-uniformity (asymptotic freedom)}

The finite gap is dimensionless. The physical gap is
$\Delta_{\mathrm{phys}}(a)=\Delta_{\mathrm{dimless}}/a$
(\verb|z2HeatTimePhysicalGap|), so naively the gap diverges as $a\to0$ unless the
dimensionless gap shrinks on a tuned schedule. Lean already exhibits a schedule
where the physical gap \emph{closes}:
\verb|z2QuadraticHeatTime_physicalGap_closes| (with $t(a)=a^2$, for every
$\varepsilon>0$ there is $a$ with $\Delta_{\mathrm{phys}}<\varepsilon$). The
crux is the opposite-direction, much harder statement:

\begin{quote}
prove a \emph{uniform positive lower bound} on the physical gap along the
continuum limit of the non-abelian theory, where asymptotic freedom forces the
coupling to run and the strong-coupling expansion that gave the finite gap is no
longer valid.
\end{quote}

\noindent This is the asymptotic-freedom wall: strong coupling (large $a$) gives
the gap easily; the continuum lives at weak coupling (small $a$), reached only
through the RG flow, and the strong-coupling control does not extend there. The
RG-bootstrap audit layer (\verb|RGBootstrap.lean|,
\verb|AffineRGStepBound.bound_by_invariant_interval|,
\verb|HasExtendedExtractionContraction|) isolates exactly which arithmetic and
contraction facts any RG proof must instantiate, and the same-constant
lower-extraction crux (\verb|not_sameConstantLowerExtractionRGCertificate| and
its even-cutoff / threshold variants) already refutes one tempting repair
corridor --- keeping the advertised constant while lowering the extraction depth.
The continuum endpoint remains the open node:
\verb|yangMillsContinuumMassGapEndpointRoadmapEntry| (0\% continuum progress).

\subsection{Possible avenues still open}

\paragraph{Y1. Non-abelian finite landmark.} The $\mathbb{Z}_2$ landmark is
abelian. The cleanest next foothold is to repeat
\verb|z2StrongCoupling_latticeYangMills_massGap_landmark| for a genuinely
non-abelian finite group ($S_3$, then $SU(2)$ truncations): an explicit
strong-coupling transfer matrix, reflection positivity, and a positive gap.
This generalizes the construction without yet touching the continuum.

\paragraph{Y2. Make the RG bootstrap analytic, not just arithmetic.} The RG
layer currently proves the \emph{finite arithmetic} of contraction
(\verb|AffineRGStepBound|, the exact $2224\cdot2^{-13}<1/3$ gain). The avenue is
to instantiate \verb|AffineRGStepBound| with an \emph{actual polymer norm
recurrence} --- the one-step RG estimate with real constants --- so the
bootstrap controls a genuine gap quantity rather than a placeholder.

\paragraph{Y3. Clustering bridge from decay to gap.} The clustering audit
(\verb|SpectralGapClusteringBridge|,
\verb|not_exists_hasSpectralMassGap_of_positiveLongRangeLowerBound|) shows that a
persistent long-range correlation tail kills the gap. The avenue is to supply,
on the lattice family, a proved exponential-decay estimate for centered
correlations that is \emph{uniform} in $a$ --- this is the clustering face of
gap-uniformity.

\subsection{Conjectured workarounds for the asymptotic-freedom wall}

\begin{conjecture}[Speculative --- Gap-uniformity from a uniform RG contraction]
\label{conj:ym-rg}
There is a single contraction constant $\kappa<1$ and an invariant interval $B$
such that the polymer norm recurrence
$a_{n+1}\le\kappa a_n+\varepsilon$ holds at \emph{every} RG step from strong
toward weak coupling, with $\varepsilon\le(1-\kappa)B$. By
\verb|AffineRGStepBound.bound_by_invariant_interval| this bounds every finite
iterate, and a uniform-in-$a$ version yields a continuum gap lower bound. The
content is that asymptotic freedom makes the irrelevant couplings contract fast
enough that the gap does not collapse.
\end{conjecture}

\noindent\emph{Why it might hold.} Asymptotic freedom is a contraction
statement: the coupling flows to zero, so the dangerous operators are
\emph{irrelevant} (the \verb|irrelevantScale| $b^{3-d_{\max}}<1$ arithmetic is
exactly this). \emph{Why it might fail.} A uniform $\kappa<1$ across the whole
flow is precisely the hard analytic estimate; the same-constant crux already
shows that cheap versions of this contraction are formally refuted, so any
workable $\kappa$ must come from a genuine multiscale estimate, not a constant
reused from strong coupling.

\begin{conjecture}[Speculative --- Reflection positivity is preserved by the
limit]
\label{conj:ym-rp-limit}
The finite reflection positivity (\verb|z2StrongCouplingKernel_reflectionPositive|,
\verb|OSReflectionPositive|) passes to the continuum measure provided the
lattice family is tight, so that the continuum-measure and reflection-positivity
gate canaries clear simultaneously. Tightness $+$ reflection positivity
$+$ Euclidean covariance then deliver an OS continuum theory whose transfer
operator inherits the lattice gap as a lower bound.
\end{conjecture}

\noindent\emph{Why it might hold.} Reflection positivity is a closed condition;
limits of reflection-positive measures are reflection-positive when the limit
exists. \emph{Why it might fail.} The limit's \emph{existence} (tightness) is the
unproved part, and the canaries
(\verb|toyYangMillsContinuumMeasureTightnessWithoutLimitAudit| etc.) prove that
no single sub-interface suffices; the conjecture is responsible only if tightness
is established independently.

\subsection{Where to point creative effort next}

The single highest-leverage target is
\textbf{Conjecture~\ref{conj:ym-rg} (uniform RG contraction), entered through
avenue Y2}. Reason: the $\mathbb{Z}_2$ landmark already proves a positive gap
exists at strong coupling; the \emph{only} thing standing between that and the
continuum is gap-uniformity, and the RG bootstrap layer has already reduced
gap-uniformity to a single recurrence with an explicit invariant-interval
bootstrap lemma. The recommended push: instantiate
\verb|AffineRGStepBound| with a concrete polymer-norm recurrence carrying real
constants and search for a uniform $\kappa<1$ valid across the strong-to-weak
flow, using the same-constant crux theorems as a guard against contractions that
merely reuse the strong-coupling constant. If a uniform $\kappa$ is found, the
finite landmark upgrades toward a continuum lower bound; if every candidate
$\kappa$ fails the same-constant guard, the failure isolates the precise
multiscale estimate that asymptotic freedom must supply --- the true heart of the
problem, now named and bounded rather than diffuse.

\subsection*{PLN lens}

\begin{table}[h]
\centering
\small
\begin{tabular}{lcccc}
\toprule
Central node & STV $\langle s,c\rangle$ & ITV $[L,U]$ & Progress & Status \\
\midrule
Finite $\mathbb{Z}_2$ strong-coupling lattice gap & $\langle 0.55, 0.30\rangle$ & $[0.165,\,0.865]$ & --- & built (landmark) \\
Continuum gap-uniformity (asymptotic freedom) & $\langle 0.05, 0.16\rangle$ & $[0.008,\,0.848]$ & --- & open crux \\
OS-gate / canary audit layer & $\langle 0.95, 0.80\rangle$ & $[0.76,\,0.96]$ & --- & built \\
\midrule
\textbf{Route (through, to continuum)} & $\langle 0.05, 0.16\rangle$ & $[0.008,\,0.848]$ & $\sim$30\% & \\
\bottomrule
\end{tabular}
\caption{PLN summary for the lattice-OS-spectral-gap route. The finite
$\mathbb{Z}_2$ lattice gap node is genuinely positive
($\langle 0.55,0.30\rangle$): a strictly positive gap is machine-checked for that
finite model. The through-route to the continuum inherits the continuum crux's
low strength, because the asymptotic-freedom wall is untouched.}
\end{table}

\noindent Lattice-gap node: $\langle 0.55, 0.30\rangle$, ITV
$[0.55\cdot 0.30,\;0.55\cdot 0.30+(1-0.30)] = [0.165,\,0.865]$.
Continuum node and through-route: $\langle 0.05, 0.16\rangle$, ITV
$[0.05\cdot 0.16,\;0.05\cdot 0.16+(1-0.16)] = [0.008,\,0.848]$,
PROGRESS $\sim 30\%$. The split is the point: the finite construction is a real,
checked success worth building on, while the continuum endpoint remains a
genuine open problem guarded by asymptotic freedom. Progress is dominated by how
much of the continuum machinery is unbuilt, not by the strength of the finite
landmark.
